% !TEX root = ../bb_AiRa_Kappaleet.tex

% ö = \%c3\%b6
% ä = \%C3\%A4

\xdef\sectiontitle{\IVsectiontitle} %aiheuttama
\TOCrefs

%%%%%%%%%%%%%%%
\begin{frame}[label=4_5_a]
\frametitle{\texorpdfstring{\culine[\sectiontitleulinecolor]{\sectiontitle}}{\sectiontitle} }
\label{\TOCname}

\vspace{-3mm}

%\vspace{\chapterTOCvksip}%
\begin{itemize}[leftmargin=0.5cm,itemsep=\chapterTOCitemsep]
    \item[4.1]<1-> \hyperlink{4_1_a}{\IVaSubsectiontitle}
    \item[4.2]<1-> \hyperlink{4_2_a}{\IVbSubsectiontitle}	
    \item[4.3]<1-> \hyperlink{4_3_a}{\IVcSubsectiontitle}	
    \item[4.4]<1-> \hyperlink{4_4_a}{\IVdSubsectiontitle}
    \item[4.5]<1-> \hyperlink{4_5_a}{\CDAlert[blue]{\IVeSubsectiontitle}}
    \item[4.6]<1-> \hyperlink{4_6_a}{\IVfSubsectiontitle}
    \item[4.7]<1-> \hyperlink{4_7_a}{\IVgSubsectiontitle}
\end{itemize}

%\endofslide 
\end{frame}
%%%%%%%%%%%%%%%

\xdef\sectiontitle{\IVeSubsectiontitle} 


%%%%%%%%%%%%%%%
\begin{frame}
\frametitle{\texorpdfstring{\culine[\sectiontitleulinecolor]{\sectiontitle}}{\sectiontitle} }
%
\framesubtitle{Tutkiminen}%Studying decay processes
\appear{}

\begin{itemize}
	\item \appear{Lähes kaikki löydetyt hiukkaset ovat epävakaita} \appear{(paitsi protoni}\appear{, elektroni} \appear{\& neutriino)}
	\item Ne hajoavat \appear{ja niiden hajoamisesta opitaan:}
	\appear{esim.~energiat}\appear{, massat}\appear{, kvanttiluvut}
	\item Useat hiukkaset ovat niin epävakaita ja lyhytikäisiä\appear{, että niitä voidaan \CDAlert[fuchsia]{tutkia vain epäsuorasti}}
	\item \CDAlert[blue]{Liikemäärä ja relativistinen kokonaisenergia säilyvät aina}:
	\appear{jos energia ei näytä säilyvän}\appear{, niin massa muuttuu}
	\item \CDAlert[red]{Varaus säilyy aina}
	\item \CDAlert[fuchsia]{Onko taustalla olemassa jotain perustavanlaatuisia lakeja?}
\end{itemize}


\endofslide 
%\endofsection
\end{frame}
%%%%%%%%%%%%%%%

%%%%%%%%%%%%%%%
\begin{frame}
\frametitle{\texorpdfstring{\culine[\sectiontitleulinecolor]{\sectiontitle}}{\sectiontitle} }
%
\framesubtitle{Tutkiminen}%Studying decay processes
\appear{}

\begin{itemize}
	\item \appear{\CDAlert[blue]{totalitarismin periaate}} \appear{$\oldrightarrow$ säilymislain perusolemus:}
	\item[] "Kaikki mikä ei ole kiellettyä, on pakollista" ~–Gell-Mann (\& R. Heinlein)%Anything that can happen, does happen"
	\item Mikäli hajoamisprosessia ei tapahdu\appear{, tälle täytyy olla syy} \implies{ On olemassa joku \CDAlert[blue]{säilymislaki}, mikä estää prosessin }
	
	
	\item Luonnossa säilymislait liittyvät \CDAlert[fuchsia]{symmetrioihin}
	
	\item Esim.\ jos hiukkasen potentiaalifunktio on vakio etenemissuunnan suhteen\appear{, potentiaali ei kohdista hiukkaseen voimaa tähän suuntaan}\appear{, jolloin hiukkasen liikemäärä säilyy}
	\implies{ vakioarvoisella potentiaalilla on translaatiosymmetriaa (\CDAlert[fuchsia]{translaatioinvariantti}) }
	
	\item Mikäli mikään vääntömomentti ei vaikuta kappaleeseen\appear{, sen potentiaali on keskeissymmetrinen}\appear{, jne.}

\end{itemize}

\endofslide 
%\endofsection
\end{frame}
%%%%%%%%%%%%%%%

%%%%%%%%%%%%%%%
\begin{frame}
\frametitle{\texorpdfstring{\culine[\sectiontitleulinecolor]{\sectiontitle}}{\sectiontitle} }
%
\framesubtitle{Emmy Noether}
\appear{}

\begin{itemize}
	\item \appear{Vastaavasti \CDAlert[red]{energian säilyminen} sulkee pois mahdollisuuden hiukkasen hajoamiselle} \appear{niin että hajoamistuotteiden kokonaismassa olisi \Alert[tunired]{suurempi kuin alkuperäinen massa} }
	
	\item Jos kineettinen energia kasvaa\appear{, niin massan täytyy pienentyä}
 \implies{ On olemassa symmetria joka liittää yhteen \CDAlert[fuchsia]{ajan ja energian}}

\item \CDAlert[black!30!fuchsia]{Emmy Noether} muodosti 1918 lausekkeen fysiikassa ilmenevien konjugaattimuuttujien suhteen. \appear{Jokainen \CDAlert[black!30!fuchsia]{fysikaalinen symmetria} \CDAlert[black!30!fuchsia]{johtaa säilymislakiin}}

\item Yleisimmät \appear{(klassiset)} \appear{säilyvät suureet} \appear{vastaavine symmetrioineen ovat:} \vspace{-1mm}
$$ \small
\begin{array}{rcl}
\appear{ \text{Energia (massa)} & \qquad \leftrightarrow & \qquad\text{ajan käännön symmetria}  \\}
\appear{ \text{Liikemäärä} & \qquad \leftrightarrow &\qquad \text{translaatiosymmetria} \\}
 \appear{\text{Kulmaliikemäärä} & \qquad \leftrightarrow  & \qquad\text{pyörähdyssymmetria} \\}
\appear{ \text{Sähkövaraus} &\qquad \leftrightarrow &  \qquad \text{sähkömagnetismin mittainvarianssi}} \\
\appear{\text{\CDAlert[blue]{Minkowskin metriikka}} &\qquad \leftrightarrow &  \qquad \text{\CDAlert[blue]{erityinen suhteellisuusteoria}}} 
\end{array} 
$$

\end{itemize}

\endofslide 
%\endofsection
\end{frame}
%%%%%%%%%%%%%%%

%%%%%%%%%%%%%%%
\begin{frame}
\frametitle{\texorpdfstring{\culine[\sectiontitleulinecolor]{\sectiontitle}}{\sectiontitle} }
%
\framesubtitle{Emmy Noether}
\appear{}

\begin{itemize}[itemsep=2.3mm]
	 \item Noether yhdisti myös \CDAlert[red]{konjugaattimuuttujat} \appear{\CDAlert[red]{säilymislakeihin}:}
	\begin{itemize}
		 \item Energia–aika (($\exp (\rmi E t)$\,))
		\item Liikemäärä–siirtymä (($\exp (\rmi \vec{k} \cdot \vec{r})$\,))
		\item jne.
	\end{itemize}
	
	\item Kuitenkin jo ennen  \CDAlert[fuchsia]{Noetherin teoreemaa}\appear{, nämä säilymislait olivat enemmän tai vähemmän tunnettuja}\appear{, klassisessa fysiikassa enimmäkseen kokeellisten havaintojen pohjalta muodostettuina}
	
	\item Hiukkasfysiikassa %(EI kvanttimekaniikassa) 
	\appear{suurin osa säilymislaista on muodostettu empiirisesti}
	
	\item Usein säilymislain on havaittu olevan voimassa vain tietyn tyyppisessä hajoamisessa 
	
	\item Esimerkiksi \CDAlert[blue]{outous} \Alert[blue]{säilyy vain vahvoissa ja sähkömagneettisissa vuorovaikutuksissa}, muttei heikoissa vuorovaikutuksissa, jne.
\end{itemize}

\endofslide 
%\endofsection
\end{frame}
%%%%%%%%%%%%%%%


%%%%%%%%%%%%%%%
\begin{frame}
\frametitle{\texorpdfstring{\culine[\sectiontitleulinecolor]{\sectiontitle}}{\sectiontitle} }
%
\framesubtitle{Emmy Noether}
\appear{}

\begin{itemize}
	\item \appear{\CDAlert[blue]{Kvanttimekaniikan säilymislait} liittyvät tiettyihin \CDAlert[blue]{kvanttilukuihin}}%
	\appear{\xdef\loc{south east}%
\begin{tikzpicture}[remember picture,overlay] % anchor=south
    \node (A) [yshift=-0.35*\figYshift,xshift=\figXshift, anchor=\loc] at (current page.\loc) {\includegraphics[page=1,width=13cm]{Figures_booklet/SOM_11_Fundamental_particles_figures/SOM_Fundamental_Table_4.png}};%scale=\figscale. % 8cm max hei
\end{tikzpicture}\CR[ref= h]}%
	
	\item Eli uusia \CDAlert[tuniblue]{kvanttilukuja tarvitaan}, jotka säilyvät tietyissä reaktioissa, prosesseissa tai vuorovaikutuksissa: \vspace{1.5mm}
	\begin{itemize} \small
	 \item \CDAlert[red]{Baryoniluku (B)}: \appear{+1 baryonille}\appear{, $-1$ antibaryonille}\appear{, 0 muulloin}
	\item \CDAlert[blue]{Leptoniluvut (L)}: \appear{yksi jokaiselle leptonisukupolvelle} \appear{$L_{e}$}\appear{, $L_{\mu}$}\appear{, $L_{\tau}$} \appear{(uutta dataa 1998)}
	\item \CDAlert[fuchsia]{Outous (S)}: \appear{$-1$ oudolle kvarkille, $+1$ antioudolle, 0 muulloin}
	\item \CDAlert[green]{Isospin}
	\item \CDAlert[black!20!antiblue]{Hypervaraus}
	\end{itemize}

\end{itemize}




\endofslide 
%\endofsection
\end{frame}
%%%%%%%%%%%%%%%


%%%%%%%%%%%%%%%
\begin{frame}
\frametitle{\texorpdfstring{\culine[\sectiontitleulinecolor]{\sectiontitle}}{\sectiontitle} }
%
\framesubtitle{Leptoniluku}%On the lepton number
\appear{}

\seminormal
\begin{itemize}
	\item \appear{Alkuperäisen standarimallin mukaan}\appear{, \CDAlert[blue]{neutriinoilla} ei pitäisi olla massaa, mutta niillä on polarisaatio} 
	\item Kokeiden perusteella neutriinot ovat vasenkätisiä\appear{, eli niiden spin on vastakkaissuuntainen niiden liikemäärän kanssa}
	\item Antineutriinot ovat oikeakätisiä\appear{, jolloin spin ja liikemäärä ovat samansuuntaiset}

\item Standardimallissa\appear{, massa on seurausta  \CDAlert[fuchsia]{Higgsin  bosonista}.} \appear{Saman vuorovaikutuksen oletetaan myös muuttavan kätisyyden oikeasta vasemmaksi, ja päinvastoin}

\item Koska on havaittu vain vasenkätisiä neutriinoita\appear{, niiden pitäisi olla \CDAlert[tuniblue]{massattomia}.} \appear{Tästä syystä leptoniluku näyttäisi säilyvän heikossa vuorovaikutuksessa}\appear{, ja tarkemmin ottaen jokainen niiden \Alert[red]{mauista säilyy}}

\item Mutta kokeet 1998  \appear{osoittivat että \Alert[blue]{neutriinot voivat muuttua toisiksi neutriinoiksi}}\appear{, ilmiössä jota kutsutaan \CDAlert[blue]{neutriino-oskillaatioksi}}

\end{itemize}

\endofslide 
%\endofsection
\end{frame}
%%%%%%%%%%%%%%%


%%%%%%%%%%%%%%%
\begin{frame}
\frametitle{\texorpdfstring{\culine[\sectiontitleulinecolor]{\sectiontitle}}{\sectiontitle} }
%
\framesubtitle{Feynmanin kaaviot}
\appear{}

\begin{itemize}
	 \item \appear{Alkeishiukkasten törmäyksiä/reaktioita kuvataan havainnollistavilla \CDAlert[red]{Feynmanin kaavioilla}}

	\item Omaksumme alla Harrisin kirjan notaation:
	\begin{itemize}
		 \item Horisontaalinen akseli kuvaa aikaa 
		\item Vertikaalinen akseli paikkaa 
	\end{itemize}
	
	\item Absoluuttiset paikat/ajat eivät merkitse kaavioissa, ne vain helpottavat \CDAlert[red]{mahdollisten vuorovaikutusten kirjanpitoa} 
	
\end{itemize}

% 6.5 cm = midway, 10 cm = 3/4 
\vspace{2.5mm}
\begin{minipage}{9cm}
\begin{itemize}
	\item \CDAlert[blue]{Fermionit}\appear{, (leptonit}\appear{, hadronit}\appear{, kvarkit)} \appear{kuvataan  \CDAlert[blue]{kiinteällä suoralla viivalla}}\appear{, \CDAlert[fuchsia]{bosonit}} \appear{(välittäjähiukkaset)} \appear{\CDAlert[fuchsia]{aaltomaisella viivalla}.} \appear{Viivojen risteyskohta tarkoittaa reaktiota}
\end{itemize}
\end{minipage}
%\vspace{2.5mm}


\appear{\xdef\loc{south east}
\begin{tikzpicture}[remember picture,overlay] % anchor=south
    \node (A) [yshift=-0.35*\figYshift,xshift=\figXshift, anchor=\loc] at (current page.\loc) {\includegraphics[page=1,width=4.5cm]{Figures_booklet/SOM_11_Fundamental_particles_figures/SOM_Fundamental_particles_Figure_10.png}}; %scale=\figscale. % 8cm max hei
\end{tikzpicture}\CR[]{}}

\endofslide 
%\endofsection
\end{frame}
%%%%%%%%%%%%%%%

%%%%%%%%%%%%%%%
\begin{frame}
\frametitle{\texorpdfstring{\culine[\sectiontitleulinecolor]{\sectiontitle}}{\sectiontitle} }
%
\framesubtitle{Feynmanin kaaviot}
\appear{}

\seminormal
\begin{itemize}
	\item Vuorovaikutusten perussäännöt: \vspace{0.3mm}
	\begin{itemize}[itemsep=1.15mm] \small
	 \item Vain hiukkaset joilla on \CDAlert[fuchsia]{väri} voivat osallistua \CDAlert[fuchsia]{vahvaan voimaan}
	\appear{$\oldrightarrow$ \Alert[red]{Vahvan voiman solmukohtiin} liittyvät fermionit ovat \CDAlert[red]{kvarkkeja}}
	\item Värin säilymisen takia \appear{gluonit ovat jossain 8:sta \CDAlert[blue]{väri–antiväri väritilasta}.} \appear{Kvarkit vaihtavat väriä!} \appear{(tilojen määrä 8 on seurausta SU(3) symmetriaryhmästä$^*$...)\footnoteleft{\appear{\tiny{$^*$Jos kiinnostaa, tutustu David Hestenesiin ja geometriseen algebraan.}}}}
	\item Kaikki hiukkaset \appear{\CDAlert[fuchsia]{paitsi neutriinot}} \appear{\CDAlert[fuchsia]{vuorovaikuttavat sähkömagneettisesti}}
	\item Heikon voiman solmukohdat näyttävät että kaikki alkeisfermionit voivat vuorovaikuttaa W ja Z bosonien avulla
	\item \CDAlert[blue]{W tapauksessa}\appear{\Alert[blue]{, kvarkin varaus muuttuu ja kvarkki} \CDAlert[blue]{transmutatoituu}}
\end{itemize}

\end{itemize}


\appear{\xdef\loc{south}
\begin{tikzpicture}[remember picture,overlay] % anchor=south
    \node (A) [yshift=-0.25*\figYshift,xshift=0*\figXshift, anchor=\loc] at (current page.\loc) {\includegraphics[page=1,width=14cm]{Figures_booklet/SOM_11_Fundamental_particles_figures/SOM_Fundamental_particles_Figure_8.png}}; %scale=\figscale. % 8cm max hei
\end{tikzpicture}\CR[]{}}

\endofslide 
%\endofsection
\end{frame}
%%%%%%%%%%%%%%%

%%%%%%%%%%%%%%%
\begin{frame}
\frametitle{\texorpdfstring{\culine[\sectiontitleulinecolor]{\sectiontitle}}{\sectiontitle} }
%
\framesubtitle{Feynmanin kaaviot}
\appear{}

\begin{itemize}
	\item \CDAlert[red]{Fermionit} \appear{(kvarkit tai leptonit)} \appear{voivat esim.\ hajota \CDAlert[red]{bosoniksi}} \appear{tai syntyä bosonista}\appear{, tätä kuvaa \CDAlert[red]{hajoamis/syntymisvuorovaikutuspisteet} (solmut)}
	\appear{\xdef\loc{south west}%
\begin{tikzpicture}[remember picture,overlay] % anchor=south
    \node (A) [yshift=-0.25*\figYshift,xshift=-\figXshift, anchor=\loc] at (current page.\loc) {\includegraphics[page=1,width=8.7cm]{Figures_booklet/SOM_11_Fundamental_particles_figures/SOM_Fundamental_particles_Figure_9.png}};%scale=\figscale. % 8cm max hei
\end{tikzpicture}\CR[ref= h]{}}%
	\item Huomaa \CDAlert[blue]{ajankäännön symmetria} kahdessa kuvaajassa kvarkeilla ja leptoneilla (kuva 9) 

\item Koko vuorovaikutus voidaan kuvata  \appear{yhdistämällä solmukohdat monimutkaisemmaksi kaavioksi}\appear{ (kuva 10)}
\end{itemize}

\appear{\xdef\loc{south east}%
\begin{tikzpicture}[remember picture,overlay] % anchor=south
    \node (A) [yshift=-0.4*\figYshift,xshift=\figXshift, anchor=\loc] at (current page.\loc) {\includegraphics[page=1,width=4.3cm]{Figures_booklet/SOM_11_Fundamental_particles_figures/SOM_Fundamental_particles_Figure_10.png}};%scale=\figscale. % 8cm max hei
\end{tikzpicture}}

\endofslide 
%\endofsection
\end{frame}
%%%%%%%%%%%%%%%


%%%%%%%%%%%%%%%
\begin{frame}
\frametitle{\texorpdfstring{\culine[\sectiontitleulinecolor]{\sectiontitle}}{\sectiontitle} }
%
\framesubtitle{Feynmanin kaaviot}
\appear{}

\begin{itemize}
	 \item<.->[] \appear{\xdef\loc{south}%
\begin{tikzpicture}[remember picture,overlay] % anchor=south
    \node (A) [yshift=-0.25*\figYshift,xshift=0*\figXshift, anchor=\loc] at (current page.\loc) {\includegraphics[page=1,width=9.5cm]{Figures_booklet/SOM_11_Fundamental_particles_figures/SOM_Fundamental_particles_Figure_11a.png}};%scale=\figscale. % 8cm max hei
\end{tikzpicture}\CR[]{}}%
\appear{\textbf{a)} Neutroni ($n$) muuntuu protoniksi ($p$)} \appear{ja elektroniksi ($e^-$)}
\end{itemize}


\xdef\loc{north}
\begin{tikzpicture}[remember picture,overlay] % anchor=south
    \node (A) [yshift=-1.5cm,xshift=0*\figXshift, anchor=\loc] at (current page.\loc) {\includegraphics[page=1,width=10cm]{Figures_booklet/SOM_11_Fundamental_particles_figures/SOM_Fundamental_particles_Figure_11t.png}}; %scale=\figscale. % 8cm max hei
\end{tikzpicture}

\endofslide 
%\endofsection
\end{frame}
%%%%%%%%%%%%%%%


%%%%%%%%%%%%%%%
\begin{frame}
\frametitle{\texorpdfstring{\culine[\sectiontitleulinecolor]{\sectiontitle}}{\sectiontitle} }
%
\framesubtitle{Feynmanin kaaviot}
\appear{}

\begin{itemize}
	 \item<1->[]\appear{\xdef\loc{north}
\begin{tikzpicture}[remember picture,overlay] % anchor=south
    \node (A) [yshift=-2.7cm,xshift=3.4cm, anchor=\loc] at (current page.\loc) {\includegraphics[page=1,width=1.75cm]{Figures_booklet/SOM_11_Fundamental_particles_figures/SOM_Fundamental_particles_Figure_11tt.png}}; 
\end{tikzpicture}\CR[]{}%
\xdef\loc{south}%
\begin{tikzpicture}[remember picture,overlay] % anchor=south
    \node (A) [yshift=-0.25*\figYshift,xshift=0*\figXshift, anchor=\loc] at (current page.\loc) {\includegraphics[page=1,width=9.5cm]{Figures_booklet/SOM_11_Fundamental_particles_figures/SOM_Fundamental_particles_Figure_11b.png}};% 
\end{tikzpicture}}% 
\appear{\textbf{b)} Pioni ($\pi^-$) hajoaa myoniksi ($\mu^-$).} \appear{Tämä tunnettu hajoaminen tapahtuu}\appear{ maan ilmakehän ulommissa kerroksissa, missä kosmiset pionit muuntuvat myoneiksi}\appear{, jotka saapuvat Maahan lyhyen ajan $2,2 \;\mu$s kuluttua} \appear{(\CDAlert[blue]{Lorentzin kontraktio})}
\end{itemize}

\endofslide 
%\endofsection
\end{frame}
%%%%%%%%%%%%%%%

\xdef\yshiftaus{4mm}

%%%%%%%%%%%%%%%
\begin{frame}
\frametitle{\texorpdfstring{\culine[\sectiontitleulinecolor]{\sectiontitle}}{\sectiontitle} }
%
\framesubtitle{Feynmanin kaaviot}
\appear{}

\begin{itemize}
	 \item<1->[]\appear{\xdef\loc{north}%
\begin{tikzpicture}[remember picture,overlay] % anchor=south
    \node (A) [yshift=-1.5cm,xshift=0*\figXshift, anchor=\loc] at (current page.\loc) {\includegraphics[page=1,width=10cm]{Figures_booklet/SOM_11_Fundamental_particles_figures/SOM_Fundamental_particles_Figure_11t.png}}; %scale=\figscale. % 8cm max hei
\end{tikzpicture}\CR[]{}%
\xdef\loc{south}%
\begin{tikzpicture}[remember picture,overlay] % anchor=south
    \node (A) [yshift=-0.25*\figYshift,xshift=0*\figXshift, anchor=\loc] at (current page.\loc) {\includegraphics[page=1,width=9.5cm]{Figures_booklet/SOM_11_Fundamental_particles_figures/SOM_Fundamental_particles_Figure_11c.png}}; %scale=\figscale. % 8cm max hei
\end{tikzpicture}}%
\vspace{\yshiftaus}%
\appear{\textbf{c)} Sigma-baryoni ($\sum^+$) hajoaa neutroniksi ($n$).} \appear{Outous ei säily!}
\end{itemize}

\endofslide 
%\endofsection
\end{frame}
%%%%%%%%%%%%%%%


%%%%%%%%%%%%%%%
\begin{frame}
\frametitle{\texorpdfstring{\culine[\sectiontitleulinecolor]{\sectiontitle}}{\sectiontitle} }
%
\framesubtitle{Feynmanin kaaviot}
\appear{}

\begin{itemize}
	 \item<1->[] \appear{\xdef\loc{north}%
\begin{tikzpicture}[remember picture,overlay] % anchor=south
    \node (A) [yshift=-1.5cm,xshift=0*\figXshift, anchor=\loc] at (current page.\loc) {\includegraphics[page=1,width=10cm]{Figures_booklet/SOM_11_Fundamental_particles_figures/SOM_Fundamental_particles_Figure_11t.png}};%
\end{tikzpicture}\CR[]{}%
\xdef\loc{south}%
\begin{tikzpicture}[remember picture,overlay] % anchor=south
    \node (A) [yshift=-0.25*\figYshift,xshift=0*\figXshift, anchor=\loc] at (current page.\loc) {\includegraphics[page=1,width=9.5cm]{Figures_booklet/SOM_11_Fundamental_particles_figures/SOM_Fundamental_particles_Figure_11d.png}};%
\end{tikzpicture}}%
\vspace{\yshiftaus}%
\appear{\textbf{d)} Myoni ($\mu^+$) hajoaa positroniksi ($e^+$)}
\end{itemize}


\endofslide 
%\endofsection
\end{frame}
%%%%%%%%%%%%%%%


%%%%%%%%%%%%%%%
\begin{frame}
\frametitle{\texorpdfstring{\culine[\sectiontitleulinecolor]{\sectiontitle}}{\sectiontitle} }
%
\framesubtitle{Feynmanin kaaviot}
\appear{}

\begin{itemize}
	 \item<1->[] \appear{\xdef\loc{north}%
\begin{tikzpicture}[remember picture,overlay] % anchor=south
    \node (A) [yshift=-1.5cm,xshift=0*\figXshift, anchor=\loc] at (current page.\loc) {\includegraphics[page=1,width=10cm]{Figures_booklet/SOM_11_Fundamental_particles_figures/SOM_Fundamental_particles_Figure_11t.png}};%
\end{tikzpicture}\CR[]{}%
\xdef\loc{south}%
\begin{tikzpicture}[remember picture,overlay] % anchor=south
    \node (A) [yshift=-0.25*\figYshift,xshift=0*\figXshift, anchor=\loc] at (current page.\loc) {\includegraphics[page=1,width=9.5cm]{Figures_booklet/SOM_11_Fundamental_particles_figures/SOM_Fundamental_particles_Figure_11e.png}};%
\end{tikzpicture}}%
\vspace{\yshiftaus}%
\appear{\textbf{e)} Sigma-baryoni ($\sum^0$) hajoaa lambda-baryoniksi ($\Lambda^0$)}
\end{itemize}

\endofslide 
%\endofsection
\end{frame}
%%%%%%%%%%%%%%%


%%%%%%%%%%%%%%%
\begin{frame}
\frametitle{\texorpdfstring{\culine[\sectiontitleulinecolor]{\sectiontitle}}{\sectiontitle} }
%
\framesubtitle{Feynmanin kaaviot}
\appear{}

\begin{itemize}
	 \item<1->[] \appear{\xdef\loc{north}%
\begin{tikzpicture}[remember picture,overlay] % anchor=south
    \node (A) [yshift=-1.5cm,xshift=0*\figXshift, anchor=\loc] at (current page.\loc) {\includegraphics[page=1,width=10cm]{Figures_booklet/SOM_11_Fundamental_particles_figures/SOM_Fundamental_particles_Figure_11t.png}}; %scale=\figscale. % 8cm max hei
\end{tikzpicture}\CR[]{}%
\xdef\loc{south}%
\begin{tikzpicture}[remember picture,overlay] % anchor=south
    \node (A) [yshift=-0.25*\figYshift,xshift=0*\figXshift, anchor=\loc] at (current page.\loc) {\includegraphics[page=1,width=9.5cm]{Figures_booklet/SOM_11_Fundamental_particles_figures/SOM_Fundamental_particles_Figure_11f.png}}; %scale=\figscale. % 8cm max hei
\end{tikzpicture}%
\appear{\xdef\loc{south}%
\begin{tikzpicture}[remember picture,overlay] % anchor=south
    \node (A) [yshift=1.5cm,xshift=4.2cm, anchor=west] at (current page.\loc) {$\leftarrow$ oikeasti 0};%scale=\figscale. % 8cm max hei
\end{tikzpicture}}}%
\vspace{\yshiftaus}%
\appear{\textbf{f)} Kaoni hajoaa toiseksi kaoniksi ja pioniksi}
\end{itemize}



\endofslide 
%\endofsection
\end{frame}
%%%%%%%%%%%%%%%


%%%%%%%%%%%%%%%
\begin{frame}
\frametitle{\texorpdfstring{\culine[\sectiontitleulinecolor]{\sectiontitle}}{\sectiontitle} }
%
\framesubtitle{Esimerkki I}
\appear{}%
\begin{itemize}
	 \item Tarkastellaan varauksen säilymistä\appear{, energia/massan}\appear{, kulmaliikemäärää/spiniä}\appear{, outoutta} \appear{ja baryoni- ja leptonilukuja}\appear{, seuraavissa kolmessa hajoamisessa.} \appear{Selvitetään mitkä ehdotetuista hajoamisista voivat tapahtua.}

$$
\begin{aligned}
\appear{\textbf{I:} \quad && \tau^{-} \quad &\oldrightarrow \quad \mu^{-} ~+~\nu_{\tau} ~+~\overline{\nu_{\mu}}  &&}\\
&&  &  &&\\
\appear{\textbf{II:}\quad && \Delta^{+} \quad &\oldrightarrow \quad n~+~K^{+}  &&\\
&&  &  &&}\\
\appear{\textbf{III:}\quad && \overline{\Lambda^{o}} \quad &\oldrightarrow \quad \bar{p}~+~\pi^{+}  }&&
\end{aligned}
$$
\end{itemize}

%$\tau^{-} \rightarrow \mu^{-}+v_{\tau}+\overline{v_{\mu}}$
%$\Xi^{o} \rightarrow K^{+}+K^{-}$
%$\Delta^{+} \rightarrow n+K^{+}$
%$\overline{\Lambda^{o}} \rightarrow \bar{p}+\pi^{+}$

\endofslide 
%\endofsection
\end{frame}
%%%%%%%%%%%%%%%


%%%%%%%%%%%%%%%
\begin{frame}
\frametitle{\texorpdfstring{\culine[\sectiontitleulinecolor]{\sectiontitle}}{\sectiontitle} }
%
\framesubtitle{Esimerkki II} %\& b
\appear{}

\appear{\xdef\loc{south east}%
\begin{tikzpicture}[remember picture,overlay] % anchor=south
    \node (A) [yshift=-0.35*\figYshift,xshift=\figXshift, anchor=\loc] at (current page.\loc) {\includegraphics[page=1,width=4.5cm]{Figures_booklet/SOM_11_Fundamental_particles_figures/SOM_Fundamental_particles_Figure_12a.png}};%
\end{tikzpicture}\CR[]{}%
}
$$
\begin{aligned}
\textbf{I:} \quad && \appear{\tau^{-}} \quad &\rightarrow \quad \appear{\mu^{-}} ~ \plus~\appear{\nu_{\tau}} ~ \plus ~\appear{\overline{\nu_{\mu}}}  &&\\
&&  &  &&\\
\appear{\text {Energia:}} \quad && \appear{1780 \mathrm{\;MeV}} \quad&\rightarrow \quad \appear{106 \mathrm{\;MeV}}  &&\quad\appear{\text {~OK} }\\
\appear{\text {Spin:}} \quad&& \appear{1 / 2}  \quad&\rightarrow \quad  \appear{1 / 2} ~ \plus ~ \appear{1 / 2} ~ \plus ~ \appear{1 / 2} &&\quad\appear{\text {{\sim}OK}}
\end{aligned}
$$
\implies{Koska myös leptoniluku säilyy\appear{,  \CDAlert[red]{hajoaminen on sallittu}} }

\endofslide 
%\endofsection
\end{frame}
%%%%%%%%%%%%%%%


%%%%%%%%%%%%%%%
\begin{frame}
\frametitle{\texorpdfstring{\culine[\sectiontitleulinecolor]{\sectiontitle}}{\sectiontitle} }
%
\framesubtitle{Esimerkki III}
\appear{}

$$
\begin{aligned}
\appear{\textbf{II:} \quad && \appear{\Delta^{+}}} \quad &\rightarrow \quad \appear{n}\appear{~+~K^{+}}  &&\\
&&  &  &&\\
\appear{\text {Energia:}} \quad && \appear{1232 \mathrm{\;MeV}} \quad&\rightarrow \quad \appear{940 ~+~494 \mathrm{\;MeV}}  &&\quad \appear{\text {EI OK}} \\
\appear{\text {Baryoniluku:}} \quad&& \appear{+1}  \quad&\rightarrow\quad \appear{+1  ~+~ 0} &&\quad\appear{\text {OK}} \\
\appear{\text {Spin:}} \quad&& \appear{3 / 2}  \quad&\rightarrow \quad  \appear{1 / 2 ~+~ 0} &&\quad \appear{\text{{\sim}OK}}
\end{aligned}
$$
\implies{Lopputuotteiden massa on suurempi kuin alkutuotteiden\appear{,  \CDAlert[blue]{hajoaminen EI OLE sallittu}} }

\endofslide 
%\endofsection
\end{frame}
%%%%%%%%%%%%%%%


%%%%%%%%%%%%%%%
\begin{frame}
\frametitle{\texorpdfstring{\culine[\sectiontitleulinecolor]{\sectiontitle}}{\sectiontitle} }
%
\framesubtitle{Esimerkki III}
\appear{}
\appear{\xdef\loc{south east}
\begin{tikzpicture}[remember picture,overlay] % anchor=south
    \node (A) [yshift=-0.45*\figYshift,xshift=\figXshift, anchor=\loc] at (current page.\loc) {\includegraphics[page=1,width=4.5cm]{Figures_booklet/SOM_11_Fundamental_particles_figures/SOM_Fundamental_particles_Figure_12b.png}}; %scale=\figscale. % 8cm max hei
\end{tikzpicture}\CR[]{}}%
$$
\begin{aligned}
\textbf{III:} \quad && \appear{\overline{\Lambda^{o}}} \quad &\rightarrow \quad \appear{\bar{p}~+~\pi^{+}}  &&\\
&&  &  &&\\
\appear{\text {Energia:}} \quad && \appear{1116 \mathrm{\;MeV}} \quad&\rightarrow \quad \appear{938 ~+~140 \mathrm{\;MeV}}  &&\quad\appear{\text {OK }} \\
\appear{\text {Baryoniluku:}} \quad&& \appear{-1}  \quad&\rightarrow\quad \appear{-1  ~+~ 0} &&\quad\appear{\text {OK }} \\
\appear{\text {Spin:}} \quad&& \appear{1 / 2 } \quad&\rightarrow \quad  \appear{1 / 2 ~+~ 0} &&\quad\appear{\text {OK }}
\end{aligned}
$$
\implies{ Kaikki ominaisuudet säilyvät\appear{, \CDAlert[red]{hajoaminen on sallittu}} }

%\endofslide 
\endofsection
\end{frame}
%%%%%%%%%%%%%%%
